% ============================================================================
% Fichier  : partie3/P03_C06_preparer_mission.tex
% Titre    : La Mission : Principes, Préparation et Décomposition
% Auteur   : MINKA MI NGUIDJOÏ Thierry Emmanuel
% Version  : 0.2 -- novembre 2026 (réécriture étoffée, alignée sur INF3226
%            C07 à C09, A0 et A1)
% ============================================================================
% Sources : ISO 19011:2018 (principes) ; ISSAI 4000 via le manuel IDI ;
% A0 « Vers l'auditologie » (hal-05701007) : axiomes 1, 2, 5, 6, proposition
% de décomposition hiérarchique, gabarit canonique ; A1 (hal-05701270).
% Miller (1956) ; Simon (1962).
% Faits MINÉQUIR (LAB-01) : ordre de mission du 16/03/2026 ; équipe
% NGO BIYONG (chef de mission) + ESSAMA ; 4 domaines D1 Accès, D2 Fichier
% des fournisseurs, D3 Circuit du mandat, D4 Exploitation ; cotation des
% risques R1-R8 inchangée (R1 et R3 à 9) ; contrôle D1.1.2 = désactivation
% à la date d'effet (2 critères).
% ============================================================================

\chapter[La mission : principes, préparation, décomposition]{La Mission : Principes, Préparation et Décomposition}
\label{chap:preparer}
\chaptermark{Préparer la mission}

\epigraph{%
    « Une mission se gagne avant le terrain. »%
}{Adage des vérificateurs}

\niveauChapitre{Intermédiaire}{Chapitres~\ref{chap:eval-sci}
et~\ref{chap:normes}.}

\begin{boxobjectifs}
\begin{itemize}
    \item Énoncer les principes de l'audit et les relier aux axiomes de
          l'auditologie.
    \item Décrire le cycle d'une mission en cinq phases, avec les entrées et
          les sorties de chacune.
    \item Rédiger un ordre de mission et des termes de référence qui fixent
          l'objectif, le périmètre, la période et les critères.
    \item Organiser la demande de pièces en trois strates, et décider si une
          pièce absente vaut \NON{} ou \NEK{}.
    \item Conduire une prise de connaissance bornée, par couches du système
          d'information.
    \item Hiérarchiser les zones d'audit en croisant le risque et la
          connaissance que l'on en a.
    \item Décomposer le domaine audité selon le gabarit canonique à huit
          niveaux, et rédiger une fiche d'assertion par contrôle.
    \item Établir le programme de travail.
\end{itemize}
\end{boxobjectifs}

\bigskip

% ============================================================================
\section*{Pourquoi ce chapitre}
% ============================================================================

Une mission sur un système d'information dispose de peu de temps et d'un
objet immense~: des milliers d'opérations, des dizaines de profils, des
journaux de plusieurs gigaoctets. Tout tester est impossible, et ce que
l'on ne regarde pas reste inconnu. La préparation décide donc de ce que la
mission saura, et de ce qu'elle ignorera. Elle se juge à une question~:
\emph{le programme de travail porte-t-il sur les zones où le risque est
élevé et la connaissance faible~?}

Ce chapitre suit la méthode du cours d'audit des systèmes d'information de
l'auteur, qui est l'instanciation, pour le praticien, du programme de
recherche de l'auditologie~\cite{minka2026auditologie}. Il en reprend les
trois outils de préparation~: le cycle de la mission, l'évaluation croisée
du risque et de la connaissance, et le gabarit canonique de décomposition.

\medskip
\textbf{Plan du chapitre.}
\begin{enumerate}
    \item Les principes de l'audit (section~\ref{sec:c06-principes}).
    \item Le cycle de la mission (section~\ref{sec:c06-cycle}).
    \item L'ordre de mission et les termes de référence
          (section~\ref{sec:c06-ordre}).
    \item Les pièces à demander (section~\ref{sec:c06-pieces}).
    \item La prise de connaissance (section~\ref{sec:c06-connaissance}).
    \item Hiérarchiser~: risque et connaissance
          (section~\ref{sec:c06-risques}).
    \item Décomposer~: le gabarit canonique et la fiche d'assertion
          (section~\ref{sec:c06-decomposer}).
    \item Le programme de travail (section~\ref{sec:c06-programme}).
\end{enumerate}

\separateur

% ============================================================================
\section{Les principes de l'audit}
\label{sec:c06-principes}
% ============================================================================

La norme ISO~19011 énonce sept principes de l'audit~\cite{iso19011_2018}~;
les normes de l'INTOSAI et de l'IIA en portent l'équivalent. Ils ne sont
pas des vœux~: chacun commande une pratique de la mission.

\begin{tableaucours}{L{3.4cm}L{4.6cm}L{4.2cm}}{%
    \tableauHeader{Principe} &
    \tableauHeader{Ce qu'il impose} &
    \tableauHeader{Dans une mission SI}
}
Intégrité & Agir honnêtement, avec diligence & Ne jamais modifier une
donnée de l'entité \\
Restitution impartiale & Rapporter fidèlement, y compris les difficultés
& Écrire les \NEK{} restants au rapport \\
Conscience professionnelle & Exercer le soin et le jugement requis &
Vérifier l'intégrité de chaque extraction \\
Confidentialité & Protéger l'information obtenue & Protéger les fichiers
de paie et de fournisseurs \\
Indépendance & Fonder l'objectivité des conclusions & Ne pas auditer un
système que l'on a conçu \\
Approche fondée sur la preuve & Conclure sur des preuves vérifiables &
Partir du système, pas des classeurs \\
Approche par les risques & Porter l'effort là où il compte & Croiser
risque et connaissance (section~\ref{sec:c06-risques}) \\
\end{tableaucours}
\vspace{-0.8em}
\begin{center}{\small\itshape Tableau~6.1. Les sept principes de l'audit
(d'après ISO~19011:2018)}\end{center}

\begin{boxinfo}
\textbf{Des principes aux axiomes.} L'article \emph{Vers l'auditologie}
montre que plusieurs de ces principes se déduisent d'axiomes plus
fondamentaux~\cite{minka2026auditologie}. L'approche fondée sur la preuve
découle de l'\emph{observabilité partielle} (axiome~1)~: l'auditeur ne voit
qu'une partie du système, il doit donc dire sur quoi il s'appuie.
L'indépendance n'est pas un axiome, mais une conséquence de
l'\emph{interaction stratégique} (axiome~4) et de l'\emph{opposabilité}
(axiome~6)~: une conclusion qu'un tiers peut contester doit venir de
quelqu'un qui n'a pas intérêt au résultat.
\end{boxinfo}

\subsection{Les principes en conflit}

Les principes se heurtent parfois. La confidentialité demande de limiter la
diffusion d'une extraction~; la restitution impartiale demande d'en
montrer les résultats. La règle du cours est de ne jamais sacrifier
l'opposabilité~: on peut anonymiser une liste d'agents dans le rapport, à
condition de conserver au dossier la liste nominative et son empreinte.

\separateur

% ============================================================================
\section{Le cycle de la mission}
\label{sec:c06-cycle}
% ============================================================================

\subsection{Cinq phases}

\begin{center}
\begin{tikzpicture}[
    ph/.style = {draw=CouleurBordDef, fill=CouleurDef, rounded corners=3pt,
        text width=2.1cm, minimum height=0.9cm, align=center,
        font=\scriptsize\bfseries\color{CouleurPrimaire}},
    arrow/.style = {-{Stealth[length=4pt]}, thick, color=CouleurSecondaire}
]
\foreach \t [count=\i] in {Préparer, Planifier, Exécuter, Conclure, Suivre}
    \node[ph] (p\i) at (2.55*\i, 0) {\t};
\foreach \i [evaluate=\i as \j using int(\i+1)] in {1,2,3,4}
    \draw[arrow] (p\i) -- (p\j);
\draw[arrow, dashed] (p5.south) -- ++(0,-0.55) -| (p1.south)
    node[pos=0.25, below, font=\scriptsize\itshape, color=CouleurSecondaire]
    {le cycle suivant part de l'état laissé par le précédent};
\end{tikzpicture}
\end{center}
\vspace{-0.2em}
\begin{center}{\small\itshape Figure~6.1. Le cycle de la mission}\end{center}

La flèche de retour n'est pas un ornement. L'auditologie pose
l'\emph{individuation temporelle} (axiome~5)~: le système audité évolue,
et son état à la mission suivante dépend de la mission
précédente~\cite{minka2026auditologie}. Le suivi n'est donc pas une
formalité~: il est la première phase du cycle suivant
(Chapitre~\ref{chap:constater}).

\subsection{Entrées et sorties de chaque phase}

\begin{tableaucours}{L{2.2cm}L{4.6cm}L{5.4cm}}{%
    \tableauHeader{Phase} &
    \tableauHeader{Entrées} &
    \tableauHeader{Sorties}
}
Préparer & Mandat, rapports antérieurs, textes & Ordre de mission, termes
de référence, demande de pièces \\
Planifier & Prise de connaissance, cotation des risques & Décomposition,
fiches d'assertion, programme de travail \\
Exécuter & Programme, pièces, accès au système & Feuilles de test, cotes
\OUI{}, \NON{}, \NEK{}, preuves et empreintes \\
Conclure & Cotes et preuves & Mesures par domaine, observations, rapport
provisoire, contradictoire, rapport définitif \\
Suivre & Plan d'action & Tableau de suivi, comparaison des cycles \\
\end{tableaucours}
\vspace{-0.8em}
\begin{center}{\small\itshape Tableau~6.2. Entrées et sorties des cinq
phases}\end{center}

\subsection{L'équipe et la supervision}

Une mission du CONSUPE compte au moins un chef de mission et un
vérificateur. Le chef de mission répond du programme, de la cohérence des
cotes et du rapport~; il vise chaque feuille de test. La supervision n'est
pas une relecture finale~: elle a lieu pendant l'exécution, quand une cote
peut encore être reprise. Au MINÉQUIR, Mme~NGO~BIYONG revoit chaque soir
les feuilles de M.~ESSAMA.

\separateur

% ============================================================================
\section{L'ordre de mission et les termes de référence}
\label{sec:c06-ordre}
% ============================================================================

\subsection{L'ordre de mission}

L'ordre de mission est la première pièce du dossier. Il fonde le droit
d'accès du vérificateur (décret n°~2013/287, art.~2, al.~4) et fixe ce que
la mission doit démontrer.

\begin{tableaucours}{L{3.0cm}L{9.2cm}}{%
    \tableauHeader{Rubrique} &
    \tableauHeader{Ordre de mission du 16 mars 2026 (MINÉQUIR)}
}
Fondement & Décret n°~2013/287, art.~2, al.~3 (contrôle des SI) et
al.~4 \\
Objectif & Évaluer le contrôle interne du circuit de la dépense de la DAF
dans SIDÉPENSE, et la régularité des paiements de l'exercice \\
Périmètre & DAF du MINÉQUIR~; application SIDÉPENSE~; poste comptable
pour la seule étape du paiement \\
Période & Exercice 2025 \\
Durée & Quatre semaines, du 16~mars au 10~avril~2026 \\
Équipe & Mme~NGO~BIYONG (chef de mission), M.~ESSAMA (vérificateur) \\
Livrables & Rapport provisoire, contradictoire, rapport définitif \\
\end{tableaucours}
\vspace{-0.8em}
\begin{center}{\small\itshape Tableau~6.3. Les rubriques d'un ordre de
mission}\end{center}

\subsection{Les termes de référence : quatre questions}

Les termes de référence précisent l'ordre de mission. Ils répondent à
quatre questions, et à elles seules.

\begin{enumerate}
    \item \textbf{Que doit-on démontrer~?} L'objectif, écrit comme une
          question fermée~: «~Le contrôle interne du circuit de la dépense
          dans SIDÉPENSE est-il efficace pour l'exercice 2025~?~»
    \item \textbf{Sur quoi~?} Le périmètre~: entités, applications,
          processus inclus, et ce qui est exclu (ici, la paie et la recette).
    \item \textbf{Par rapport à quoi~?} Les critères et leur source
          (Chapitre~\ref{chap:normes}).
    \item \textbf{Avec quelle assurance~?} Raisonnable ou limitée
          (Chapitre~\ref{chap:audit-public}), et les seuils de décision que
          la mission retiendra (Chapitre~\ref{chap:constater}).
\end{enumerate}

\begin{boxalert}
\textbf{Violation typique.} Un ordre de mission qui dit «~contrôler le
système informatique de la DAF~», sans objectif, sans période et sans
périmètre.

\smallskip
\textbf{Règle~:} un objectif s'écrit comme une question à laquelle le
rapport répondra par oui ou par non~; le périmètre dit ce qui est exclu
autant que ce qui est inclus.
\end{boxalert}

\separateur

% ============================================================================
\section{Les pièces à demander}
\label{sec:c06-pieces}
% ============================================================================

\subsection{Trois strates}

La demande de pièces part dès le premier jour, par écrit, avec un délai.
Elle s'organise en trois strates, de la plus générale à la plus probante.

\begin{tableaucours}{L{2.8cm}L{4.8cm}L{4.6cm}}{%
    \tableauHeader{Strate} &
    \tableauHeader{Contenu} &
    \tableauHeader{Au MINÉQUIR}
}
1. Gouvernance & Organigramme, décisions, schéma directeur, politiques &
Organigramme de la DAF, décision de création de la cellule \\
2. Fonctionnement & Procédures, manuels, contrats, comptes rendus &
Procédures du circuit, contrat de maintenance de SIDÉPENSE \\
3. Traces & Extractions, journaux, paramètres, faites devant l'auditeur &
Comptes et profils, mandats 2025, journal des opérations \\
\end{tableaucours}
\vspace{-0.8em}
\begin{center}{\small\itshape Tableau~6.4. Les trois strates de
pièces}\end{center}

\begin{tableaucours}{cL{6.2cm}L{3.0cm}L{2.2cm}}{%
    \tableauHeader{N°} &
    \tableauHeader{Pièce} &
    \tableauHeader{Destinataire} &
    \tableauHeader{Délai}
}
1 & Liste des utilisateurs de SIDÉPENSE et de leurs profils, extraite
devant la vérificatrice & Cellule informatique & J+2 \\
2 & Liste des mutations et départs de 2025 & Ressources humaines & J+3 \\
3 & Extraction des mandats 2025 et du journal des opérations & Cellule
informatique & J+5 \\
4 & Journal des modifications du fichier des fournisseurs & Cellule
informatique & J+5 \\
5 & Procédures écrites du circuit et de la gestion des comptes & DAF &
J+3 \\
6 & Dernier rapport d'audit de sécurité de l'ANTIC (date, périmètre, plan
d'action) & Cellule informatique & J+5 \\
\end{tableaucours}
\vspace{-0.8em}
\begin{center}{\small\itshape Tableau~6.5. Demande de pièces de la mission
MINÉQUIR}\end{center}

\subsection{Une pièce absente : NON ou NEK~?}

Une pièce qui ne vient pas n'a pas toujours le même sens. Le cours
applique une règle simple, cohérente avec l'instrument de l'article
A1~\cite{minka2026nek}.

\begin{itemize}
    \item \textbf{\NON{}}~: la pièce devait exister (une règle l'impose), et
          l'entité confirme par écrit qu'elle n'existe pas. L'absence est
          une constatation.
    \item \textbf{\NEK{}}~: la pièce peut exister ailleurs, ou n'a pas
          encore été produite malgré la relance. On ne sait pas~; on le dit,
          avec les diligences faites.
\end{itemize}
Un \NEK{} qui reste à la fin de la mission devient une \emph{limitation
d'étendue}, écrite au rapport (Chapitre~\ref{chap:constater}).

\separateur

% ============================================================================
\section{La prise de connaissance}
\label{sec:c06-connaissance}
% ============================================================================

\subsection{Une prise de connaissance bornée}

La prise de connaissance n'a pas pour but de tout comprendre, mais de
savoir assez pour formuler des hypothèses de risque. L'article A0 rappelle
que l'auditeur est un agent cognitivement limité
(axiome~2)~\cite{minka2026auditologie}~: une prise de connaissance sans
borne produit un dossier que personne ne relit. Le cours la limite à deux
jours sur quatre semaines, et la conduit par couches.

\subsection{Lire le système d'information en couches}

\begin{tableaucours}{L{2.8cm}L{4.8cm}L{4.6cm}}{%
    \tableauHeader{Couche} &
    \tableauHeader{Question} &
    \tableauHeader{Au MINÉQUIR}
}
Organisation & Qui décide, qui exécute, qui contrôle~? & DAF propriétaire,
cellule exécutante, contrôle financier \\
Processus & Quelles étapes, dans quel ordre~? & Engagement, liquidation,
ordonnancement, paiement \\
Applications & Quelles applications portent les étapes~? & SIDÉPENSE,
logiciel du poste comptable \\
Données & Quelles données de référence conditionnent tout~? & Fichier des
fournisseurs, table des profils \\
Infrastructure & Où tourne le système, qui l'administre~? & Serveur de la
cellule, sauvegardes locales \\
\end{tableaucours}
\vspace{-0.8em}
\begin{center}{\small\itshape Tableau~6.6. Les couches du système
d'information}\end{center}

\subsection{La cartographie des applications et des flux}

\begin{tableaucours}{L{2.8cm}L{3.2cm}L{3.0cm}L{3.2cm}}{%
    \tableauHeader{Application} &
    \tableauHeader{Fonction} &
    \tableauHeader{Propriétaire} &
    \tableauHeader{Flux sortant}
}
SIDÉPENSE & Engagement, liquidation, ordonnancement, fournisseurs & DAF
& Mandats vers le poste comptable (fichier quotidien) \\
Logiciel du poste comptable & Paiement, comptabilité & Poste comptable &
Ordres de virement vers les banques \\
Fichier du personnel & Mutations, départs & Direction des ressources
humaines & Aucun vers SIDÉPENSE (constat) \\
\end{tableaucours}
\vspace{-0.8em}
\begin{center}{\small\itshape Tableau~6.7. Cartographie des applications
du périmètre}\end{center}

La dernière ligne est déjà un résultat~: aucun flux ne relie le fichier du
personnel à SIDÉPENSE. Les mutations ne sont connues de la cellule
informatique que si quelqu'un pense à les lui signaler. La cartographie a
désigné un risque avant tout test.

\subsection{Conduire un entretien de prise de connaissance}

\begin{boxprocedure}
\textbf{Entretien de prise de connaissance}~:
\begin{itemize}
    \item préparer une liste de questions ouvertes, par couche~;
    \item être deux~: l'un interroge, l'autre note~;
    \item demander, pour chaque affirmation, la pièce qui la montre~;
    \item relire le compte rendu à l'interlocuteur et le faire valider~;
    \item noter à part ce qui est déclaré sans pièce~: ce sont les futurs
          \NEK{} à lever.
\end{itemize}
\end{boxprocedure}

\separateur

% ============================================================================
\section{Hiérarchiser : risque et connaissance}
\label{sec:c06-risques}
% ============================================================================

\subsection{La cotation du risque}

Pour chaque risque de la matrice (Chapitre~\ref{chap:eval-sci}), on cote
de 1 (faible) à 3 (élevé) le \termecle{risque inhérent}\index{risque
inhérent} (gravité et probabilité de l'anomalie si aucun contrôle
n'existait) et le \termecle{risque lié au contrôle}\index{risque lié au
contrôle} (probabilité que le contrôle interne, tel qu'on le connaît, ne
l'empêche ni ne la détecte). Leur produit, de 1 à 9, est la priorité.

\begin{tableaucours}{L{5.4cm}ccc}{%
    \tableauHeader{Risque} &
    \tableauHeader{Inhérent} &
    \tableauHeader{Contrôle} &
    \tableauHeader{Priorité}
}
R1. Un agent parti agit encore & 3 & 3 & 9 \\
R2. Un agent émet seul un mandat & 3 & 2 & 6 \\
R3. Le bénéficiaire est changé & 3 & 3 & 9 \\
R4. Une facture est payée deux fois & 3 & 2 & 6 \\
R5. Le mandat n'est pas justifié & 2 & 1 & 2 \\
R6. Une opération ne laisse pas de trace & 2 & 1 & 2 \\
R7. Une anomalie de fin de gestion passe & 2 & 3 & 6 \\
R8. Les traces disparaissent & 2 & 2 & 4 \\
\end{tableaucours}
\vspace{-0.8em}
\begin{center}{\small\itshape Tableau~6.8. Cotation préliminaire des
risques du circuit de la dépense (convention du cours)}\end{center}

\subsection{Le risque et la connaissance}

La cotation du risque ne suffit pas. L'auditologie ajoute une seconde
dimension~: ce que l'on \emph{sait} de chaque zone. Une zone à risque élevé
que l'on connaît bien (un audit récent l'a testée) appelle moins d'effort
qu'une zone à risque moyen dont on ignore tout. C'est la distinction d'A0
entre non-conformité \emph{connue} et non-conformité
\emph{inconnue}~\cite{minka2026auditologie}~: la seconde est la plus
dangereuse, parce qu'elle ne se voit pas.

\begin{tableaucours}{L{3.0cm}L{4.4cm}L{4.8cm}}{%
    \tableauHeader{} &
    \tableauHeader{Connaissance élevée} &
    \tableauHeader{Connaissance faible}
}
\textbf{Risque élevé} & Tester ce qui a changé depuis le dernier audit &
\textbf{Priorité absolue}~: tests étendus, preuves de strate~3 \\
\textbf{Risque faible} & Cheminement ou revue analytique & Prise de
connaissance complémentaire, puis décider \\
\end{tableaucours}
\vspace{-0.8em}
\begin{center}{\small\itshape Tableau~6.9. Croiser le risque et la
connaissance}\end{center}

\begin{boxscenario}{Dossier MINÉQUIR. Où regarder d'abord}
\textbf{Constatation.} R1 (agents partis) et R3 (RIB) sont cotés 9. Aucun
audit n'a jamais porté sur SIDÉPENSE~; l'inspection de janvier n'a examiné
que des dossiers papier.

\medskip
\textbf{Analyse} (confiance moyenne). Les deux risques sont à la fois
élevés et inconnus~: ce sont les zones de priorité absolue. La
vérificatrice leur consacre la première semaine, puis les trois risques
cotés~6. Elle ne teste R5 et R6 que par un cheminement~: leur priorité ne
justifie pas un échantillon.
\end{boxscenario}

\separateur

% ============================================================================
\section{Décomposer : le gabarit canonique et la fiche d'assertion}
\label{sec:c06-decomposer}
% ============================================================================

\subsection{Pourquoi décomposer}

Un auditeur ne peut pas examiner «~les accès~» d'un bloc. George Miller a
montré que la mémoire de travail ne retient qu'environ sept éléments à la
fois, plus ou moins deux~\cite{miller1956}~; Herbert Simon, que les
systèmes complexes sont hiérarchiques et presque
décomposables~\cite{simon1962complexity}. L'article A0 en tire une
proposition~: tout instrument qui couvre un espace d'audit non trivial doit
être organisé en une décomposition hiérarchique dont chaque niveau respecte
cette borne~\cite{minka2026auditologie}.

\begin{boxdef}
\termecle{Espace d'audit}\index{espace d'audit}~: ensemble des objets sur
lesquels porte une mission (processus, contrôles, données, traces), tel que
le délimite l'ordre de mission.

\medskip
\termecle{Décomposition}\index{décomposition}~: découpage hiérarchique de
l'espace d'audit en niveaux emboîtés, chaque élément d'un niveau se divisant
en éléments du niveau inférieur.
\end{boxdef}

La décomposition apporte quatre choses~: la \textbf{couverture} (on sait ce
qui a été audité et ce qui ne l'a pas été), la \textbf{traçabilité} (chaque
conclusion remonte à une trace, ce qui la rend opposable, axiome~6), la
\textbf{localisation de l'incertitude} (les \NEK{} se situent à un endroit
précis) et la \textbf{répartition du travail}.

\subsection{Le gabarit canonique à huit niveaux}

Le programme de l'auditologie propose un \termecle{gabarit
canonique}\index{gabarit canonique} de décomposition à huit
niveaux~\cite{minka2026auditologie}. Les quatre premiers disent \emph{ce
qui est audité}~; les quatre derniers, \emph{comment on le vérifie}. Le
contrôle fait la charnière~: c'est l'objet sur lequel l'auditeur se
prononce.

\begin{center}
\begin{tikzpicture}[
    niv/.style = {draw=CouleurBordDef, fill=CouleurDef, rounded corners=3pt,
        minimum width=2.9cm, minimum height=0.58cm, align=center,
        font=\footnotesize\bfseries\color{CouleurPrimaire}},
    note/.style = {font=\scriptsize\itshape\color{CouleurSecondaire}, anchor=west, text width=7.2cm},
    fleche/.style = {-{Stealth[length=4pt]}, thick, color=CouleurSecondaire}
]
\foreach \t/\q [count=\i] in {
    1. Domaine/D1 Accès à SIDÉPENSE,
    2. Sous-domaine/D1.1 Cycle de vie des comptes,
    3. Processus/Mutation ou départ d'un agent,
    4. Activité/Retrait des accès,
    5. Contrôle/D1.1.2 Désactivation à la date d'effet,
    6. Test/Rapprocher mutations et comptes actifs,
    7. Preuve/Extraction des comptes et son empreinte,
    8. Trace/Champ « date de désactivation » et journal} {
    \node[niv] (n\i) at (0, -0.78*\i) {\t};
    \node[note] at (1.6, -0.78*\i) {\q};
}
\foreach \i [evaluate=\i as \j using int(\i+1)] in {1,...,7}
    \draw[fleche] (n\i) -- (n\j);
\end{tikzpicture}
\end{center}
\vspace{-0.2em}
\begin{center}{\small\itshape Figure~6.2. Le gabarit canonique appliqué au
contrôle des agents partis (d'après~\cite{minka2026auditologie})}\end{center}

\begin{boxalert}
\textbf{Violation typique.} Mettre dans la même liste «~une politique
d'accès existe~» (un document) et «~le journal est activé~» (une trace),
au même niveau que «~les accès~» (un domaine).

\smallskip
\textbf{Règle~:} chaque élément se place à un seul niveau, celui de sa
nature. Un mélange de niveaux rend la couverture invérifiable.
\end{boxalert}

\subsection{Preuve et trace : deux choses différentes}

La \emph{trace} existe dans le système, que l'auditeur passe ou non~: une
entrée de journal, un champ de date. La \emph{preuve} est ce que l'auditeur
en tire et conserve~: l'extraction horodatée, son empreinte, la requête
qui l'a produite. Une trace non extraite n'est pas une preuve~; une preuve
qui ne renvoie à aucune trace (une capture d'écran sans contexte) est
fragile.

\subsection{La décomposition de la mission MINÉQUIR}

\begin{center}
\begin{tikzpicture}[
    niv/.style = {draw=CouleurBordDef, fill=CouleurDef, rounded corners=3pt,
        text width=2.9cm, minimum height=0.8cm, align=center,
        font=\scriptsize\bfseries\color{CouleurPrimaire}},
    arrow/.style = {-{Stealth[length=4pt]}, thick, color=CouleurSecondaire}
]
\node[niv, text width=6cm] (obj) at (5.1,0) {Objectif~: le contrôle interne
du circuit de la dépense est-il efficace~?};
\node[niv] (d1) at (0,-1.6) {D1 Accès};
\node[niv] (d2) at (3.4,-1.6) {D2 Fichier des fournisseurs};
\node[niv] (d3) at (6.8,-1.6) {D3 Circuit du mandat};
\node[niv] (d4) at (10.2,-1.6) {D4 Exploitation};
\foreach \d in {d1,d2,d3,d4} \draw[arrow] (obj.south) -- (\d.north);
\node[niv, font=\tiny] at (0,-3.0) {Contrôles 1, 2};
\node[niv, font=\tiny] at (3.4,-3.0) {Contrôle 3};
\node[niv, font=\tiny] at (6.8,-3.0) {Contrôles 4, 5, 7};
\node[niv, font=\tiny] at (10.2,-3.0) {Contrôles 6, 8};
\end{tikzpicture}
\end{center}
\vspace{-0.4em}
\begin{center}{\small\itshape Figure~6.3. Les quatre domaines de la
mission MINÉQUIR et les huit contrôles du prologue}\end{center}

\subsection{La fiche d'assertion}

Le gabarit dit \emph{où} se trouve un contrôle. La \termecle{fiche
d'assertion}\index{fiche d'assertion} dit \emph{ce que l'auditeur affirme}
à son sujet, et ce qui arrive si l'affirmation est fausse.

\begin{tableaucours}{L{2.4cm}L{5.0cm}L{4.8cm}}{%
    \tableauHeader{Attribut} &
    \tableauHeader{Rôle} &
    \tableauHeader{Forme}
}
Objectif & Ce que le contrôle doit garantir & À l'infinitif~: « S'assurer
que… » \\
Critère & La condition vérifiable qui définit la conformité & Au présent~:
une condition unique, observable \\
Question & La question posée pour coter le contrôle & Interrogative~;
réponse \OUI{}, \NON{} ou \NEK{} \\
Risque & L'état de non-conformité si la réponse est \NON{} & L'antithèse
du critère \\
Conséquence & L'effet sur l'entité si le risque se réalise & Financier,
opérationnel, juridique \\
\end{tableaucours}
\vspace{-0.8em}
\begin{center}{\small\itshape Tableau~6.10. Les attributs de la fiche
d'assertion}\end{center}

\begin{boxobjet}{D1.1.2 · Désactivation à la date d'effet}
\textbf{Position~:} D1 Accès $\rightarrow$ D1.1 Cycle de vie des comptes
$\rightarrow$ mutation ou départ $\rightarrow$ retrait des accès
$\rightarrow$ contrôle D1.1.2.

\textbf{Objectif~:} s'assurer qu'aucun agent muté ou parti ne peut agir
dans SIDÉPENSE.

\textbf{Critère~1~:} le compte est désactivé au plus tard à la date d'effet
de la mutation. \textbf{Question~1~:} chaque compte d'agent muté en 2025
a-t-il été désactivé au plus tard à la date d'effet~?

\textbf{Critère~2~:} la DRH notifie par écrit chaque mutation à la cellule
informatique. \textbf{Question~2~:} chaque mutation de 2025 a-t-elle fait
l'objet d'une notification écrite~?

\textbf{Test~:} rapprocher la liste des 31 mutations et départs de 2025 de
la liste des comptes actifs au 23~mars~2026~; demander les notifications.

\textbf{Preuve et trace~:} extractions et empreintes~; champ de date de
désactivation~; notifications.

\textbf{Risque~:} un compte resté actif est utilisé par son ancien
titulaire ou par un tiers.

\textbf{Conséquence~:} paiement indu, impossibilité d'imputer une
opération à une personne, suites devant le CDBF.
\end{boxobjet}

\subsection{De la fiche à la cotation}

Chaque critère reçoit sa propre question et sa propre cote. La cote du
contrôle s'en déduit par une règle de composition~: un seul critère
\NON{} suffit pour que le contrôle soit \NON{}~; à défaut, un seul critère
\NEK{} donne \NEK{}~; le contrôle n'est \OUI{} que si tous ses critères le
sont.

\begin{tableaucours}{L{3.4cm}L{3.4cm}cL{3.8cm}}{%
    \tableauHeader{Critère 1 (désactivé)} &
    \tableauHeader{Critère 2 (notifié)} &
    \tableauHeader{Cote D1.1.2} &
    \tableauHeader{Lecture}
}
\OUI{} & \OUI{} & \OUI{} & Contrôle conforme \\
\OUI{} & \NON{} & \NON{} & Désactivé par chance, sans procédure \\
\NEK{} & \OUI{} & \NEK{} & Notifié, mais désactivation non établie \\
\NON{} & \NEK{} & \NON{} & Un défaut établi suffit \\
\end{tableaucours}
\vspace{-0.8em}
\begin{center}{\small\itshape Tableau~6.11. Composer les cotes de deux
critères}\end{center}

\subsection{Les règles de qualité}

\begin{itemize}
    \item \textbf{Atomicité}~: l'objectif, le critère et la question ne
          contiennent ni «~et~» ni «~ou~»~; toute assertion composée est
          scindée.
    \item \textbf{Exhaustivité}~: les éléments d'un niveau couvrent
          entièrement l'élément du niveau supérieur.
    \item \textbf{Non-redondance}~: un contrôle n'appartient qu'à un
          domaine, sinon il est compté deux fois.
    \item \textbf{Borne cognitive}~: au plus sept à neuf éléments sous un
          même parent~; au-delà, on ajoute un niveau.
\end{itemize}

\separateur

% ============================================================================
\section{Le programme de travail}
\label{sec:c06-programme}
% ============================================================================

\begin{boxprocedure}
\textbf{Le programme de travail}, une ligne par contrôle~:
\begin{itemize}
    \item le code et l'intitulé du contrôle (position dans le gabarit)~;
    \item le risque couvert, sa priorité et la connaissance que l'on en a~;
    \item la fiche d'assertion (critères, questions)~;
    \item le test prévu (cheminement, échantillon, essai, analyse de toute
          la population) et sa taille (Chapitre~\ref{chap:preuve})~;
    \item les pièces à demander, par strate~;
    \item les seuils de décision retenus (Chapitre~\ref{chap:constater})~;
    \item le membre de l'équipe et la date prévue.
\end{itemize}
\end{boxprocedure}

\begin{boxmetier}{Vérificateur}
Gardez la trace de chaque relance. Une pièce non remise à l'échéance,
malgré relance écrite, devient soit un \NON{} (si elle devait exister),
soit une limitation d'étendue à mentionner au rapport.
\end{boxmetier}

\begin{boxmetier}{Inspecteur}
Même pour une inspection courte, écrivez l'objectif et le périmètre en
tête de votre note de mission~: ils déterminent la portée de votre
conclusion (Chapitre~\ref{chap:prologue}).
\end{boxmetier}

\begin{boxmetier}{Auditeur interne}
Votre plan d'audit annuel est une cotation des risques et des
connaissances à l'échelle de l'entité~: les applications financières, et
leurs jointures, doivent y figurer, avec la date du dernier audit qui les
a testées.
\end{boxmetier}

\separateur

\begin{boxresume}
\begin{itemize}
    \item \textbf{Sept principes, des axiomes derrière.} La preuve vient de
          l'observabilité partielle~; l'indépendance, de l'interaction
          stratégique et de l'opposabilité.
    \item \textbf{Un cycle en cinq phases}~: préparer, planifier, exécuter,
          conclure, suivre~; le suivi ouvre le cycle suivant.
    \item \textbf{Quatre questions de référence}~: quoi démontrer, sur quoi,
          par rapport à quoi, avec quelle assurance.
    \item \textbf{Trois strates de pièces}~; une pièce absente est \NON{} si
          elle devait exister, \NEK{} sinon.
    \item \textbf{Risque et connaissance.} La priorité absolue va aux zones
          à risque élevé que l'on connaît mal.
    \item \textbf{Le gabarit canonique}~: huit niveaux, du domaine à la
          trace~; une fiche d'assertion par contrôle, une cote par critère.
    \item \textbf{Un programme ligne par ligne}, écrit avant le terrain.
\end{itemize}
\end{boxresume}

\begin{boxexo}[title={\ding{45}\quad Exercice~6.1. Un ordre de mission \niveaubase}]
Rédigez l'ordre de mission et les termes de référence d'une vérification
de la paie de la CAPNO (Garoua) pour l'exercice 2025, en répondant aux
quatre questions de la section~\ref{sec:c06-ordre}.
\end{boxexo}

\begin{boxexo}[title={\ding{45}\quad Exercice~6.2. Risque et connaissance \niveaumoyen}]
Pour la paie de la CAPNO, cotez au moins six risques, indiquez pour chacun
la connaissance dont vous disposez (élevée ou faible) et placez-les dans le
Tableau~6.9. Dites lesquels vous testerez par analyse de toute la
population, et lesquels par cheminement.
\end{boxexo}

\begin{boxexo}[title={\ding{45}\quad Exercice~6.3. Une fiche d'assertion \niveaumoyen}]
Rédigez la fiche d'assertion du contrôle «~seconde validation de toute
modification de RIB~» (domaine D2), avec sa position dans le gabarit, deux
critères atomiques, leurs questions, le test, la preuve, la trace, le
risque et la conséquence. Donnez la cote du contrôle si le premier critère
est \OUI{} et le second \NEK{}.
\end{boxexo}

\begin{boxexo}[title={\ding{45}\quad Exercice~6.4. Décomposer e-RECETTES \niveauexpert}]
Décomposez l'objectif «~le système e-RECETTES garantit-il l'exhaustivité
des recettes enregistrées~?~» selon le gabarit canonique jusqu'au niveau
du contrôle, pour au moins trois sous-domaines. Vérifiez les quatre règles
de qualité et signalez une violation que vous avez dû corriger.
\end{boxexo}

\begin{boxinfo}
\textbf{Transition.} Le programme dit quoi tester, et où se trouve chaque
contrôle. Le Chapitre~\ref{chap:preuve} dit comment répondre à chaque
question~: avec quelle preuve, et que faire quand on ne sait pas.
\end{boxinfo}
